\section{Panel Principal (Dashboard)}

Toda tu información y pacientes vinculados se reflejan automáticamente en el \textbf{Dashboard (Panel Principal)}. Al iniciar sesión, entrarás en contacto directo con esta lista de trabajo (Worklist), que agrupa de forma dinámica los exámenes compartidos.

\begin{enumerate}
    \item En el extremo lateral o superior, visualizarás un pequeño menú de navegación. Ahí se ubica un ícono de tu perfil con tu nombre. Al hacer clic, podrás ajustar configuraciones primarias (como subir una foto tuya para que los pacientes te ubiquen preferentemente).
    \item Justo debajo del encabezado, reposa la matrícula visual llamada \textbf{"Mis Pacientes Vinculados" o "Estudios Compartidos"}.
    \item Verás la información en forma de tarjetas, las cuales muestran el nombre del paciente, su edad (o fecha de nacimiento), tipo de estudio y un botón de estatus codificado por colores indicando su progreso.
\end{enumerate}

\begin{figure}[H]
    \centering
    \includegraphics[width=0.9\textwidth]{screenshots/associate_doctor_dashboard.png}
    \caption{El Dashboard del Médico Asociado mostrando las listas integrales.}
\end{figure}

\subsection{El Principio de Vinculación y Privacidad}

\textbf{Aviso Operativo Importante:} Al ingresar por primera vez, tu lista de pacientes y estudios aparecerá \textbf{complementamente vacía}.
El sistema no funge como un buscador abierto; por privacidad nacional del expediente médico, nosotros impedimos a los médicos ajenos buscar indiscriminadamente la información de los pacientes de la red. 

Para que los expedientes de tus pacientes abarroten tu pantalla de forma expedita, \textbf{debe ser el propio paciente quien inicie sesión en su portal y te otorgue la autorización ("Agregue tu nombre")} dentro de su menú. En cuanto el usuario oprima conceder en su sistema, inmediatamente su nombre aparecerá dentro de tu panel para tu comodidad.
